% ECONOMICS_PROBLEM: {"id": "C72-2", "title": "A polynomial-time approximation scheme for bimatrix Nash equilibrium", "jel_code": "C72", "source_id": "OP-0131"}
% !TeX program = xelatex
\documentclass[11pt,a4paper]{article}
\usepackage[margin=23mm]{geometry}
\usepackage{fontspec}
\setmainfont{Latin Modern Roman}
\usepackage{amsmath,amssymb,mathtools}
\usepackage{xcolor,enumitem,booktabs,longtable,array,xurl,hyperref}
\definecolor{muted}{HTML}{53636C}
\hypersetup{colorlinks=true,urlcolor=blue,linkcolor=blue}
\setlength{\parindent}{0pt}
\setlength{\parskip}{5pt}
\newcommand{\R}{\mathbb R}
\newcommand{\E}{\mathbb E}
\newcommand{\N}{\mathbb N}
\newcommand{\ind}{\mathbf 1}
\newcommand{\Var}{\operatorname{Var}}
\newcommand{\Cov}{\operatorname{Cov}}
\newcommand{\supp}{\operatorname{supp}}
\DeclareMathOperator*{\argmin}{arg\,min}
\DeclareMathOperator*{\argmax}{arg\,max}
\newcommand{\entrymeta}[3]{{\sffamily\small\color{muted}\textbf{\texttt{#1}}\quad|\quad #2\par #3\par}}
\newcommand{\Problemref}[1]{\texttt{#1}}
\newcommand{\JELref}[2]{\texttt{#2} (JEL \texttt{#1})}
\begin{document}
\section*{A polynomial-time approximation scheme for bimatrix Nash equilibrium}
\textbf{Problem ID:} \texttt{C72-2}\quad \textbf{JEL:} \texttt{C72}\par
% BEGIN ECONOMICS STATEMENT
\label{problem:OP-0131}
% GAME_THEORY_PROBLEM: {"local_id": "GT-001", "original_number": 1, "kind": "P", "entry_kind": "statement_only", "published": false, "review_required": true, "category": "game_theory"}
\label{game:GT-001}
{\sffamily\small\color{muted}
\textbf{GT-001} | Algorithmic equilibrium and complexity\\
\textbf{P} | Source-linked mathematical question; strong conditional negative result known.
\par}
\subsection*{Definitions and mathematical statement}
For $r\geq1$, put $\Delta_r=\{z\in\mathbb R_{\geq0}^{r}:\sum_jz_j=1\}$. A rational bimatrix game is a pair $A,B\in\mathbb Q^{m\times n}\cap[0,1]^{m\times n}$, given by explicit binary encodings of their entries. Write $L$ for the total input length. The expected payoffs at $(x,y)\in\Delta_m\times\Delta_n$ are $x^{\mathsf T}Ay$ and $x^{\mathsf T}By$. Its additive Nash residual is
\[
 r_{A,B}(x,y)=\max\left\{\max_{1\leq i\leq m}(Ay)_i-x^{\mathsf T}Ay,
 \max_{1\leq j\leq n}(x^{\mathsf T}B)_j-x^{\mathsf T}By\right\}.
\]
An $\varepsilon$-Nash equilibrium has residual at most $\varepsilon$. The PTAS question is whether there is a deterministic algorithm $\mathcal A$ and, for every fixed rational $\varepsilon\in(0,1)$, a polynomial $p_\varepsilon$ such that, for every such game, $\mathcal A(A,B,\varepsilon)$ outputs rational $x,y$ with $r_{A,B}(x,y)\leq\varepsilon$ in at most $p_\varepsilon(L)$ bit operations.
\subsection*{Short English statement}
Can an arbitrarily accurate constant-error Nash equilibrium always be computed in polynomial time?
\subsection*{Variants and scope boundaries}
The exponent of $p_\varepsilon$ may depend on $\varepsilon$. Requiring a polynomial in both $L$ and $1/\varepsilon$ is the stronger FPTAS requirement. Randomized algorithms require a specified success probability and their own hardness assumptions. Payoff residual, geometric distance to an exact equilibrium, and well-supported approximation are different measures.
\subsection*{Source relationship and status}
The unconditional algorithmic question is not settled here. It must not be described as having no negative evidence: Rubinstein's Theorem 1.2 rules out a PTAS under the Exponential Time Hypothesis for PPAD. The conjectured hypothesis, rather than an unconditional separation, is the qualification.
\subsection*{References}
\begingroup\footnotesize\raggedright
{\footnotesize\raggedright
\textbf{[S1]} Hanyu Li, Wenhan Huang, Zhijian Duan, David Henry Mguni, Kun Shao, Jun Wang, and Xiaotie Deng (2023). \emph{A survey on algorithms for Nash equilibria in finite normal-form games}. arXiv:2312.11063. Locator: Section 3, approximation algorithms; concluding ``Open problems and future directions.'' \url{https://arxiv.org/pdf/2312.11063}\par
\textbf{[S2]} Aviad Rubinstein (2016). \emph{Settling the complexity of computing approximate two-player Nash equilibria}. arXiv:1606.04550. Locator: Hypothesis 1 and Theorem 1.2. \url{https://arxiv.org/pdf/1606.04550}\par}
\par\endgroup

\subsection*{Common conventions: Source mathematical conventions}

\textbf{Finite games.} For a finite set $A$, $\Delta(A)$ is its probability
simplex. Players choose independent mixed strategies in Nash equilibrium;
correlation is allowed only when explicitly part of the solution concept.
In a game with payoff functions $u_i$ and strategy profile $x$, an additive
$\varepsilon$ Nash equilibrium satisfies
\[
 u_i(x)\geq u_i(x_i',x_{-i})-\varepsilon
 \quad\text{for every player $i$ and every admissible deviation $x_i'$}.
\]
For numerical approximation, payoffs are normalized as locally specified.
An $\varepsilon$ payoff residual is distinct from distance at most
$\varepsilon$ to the set of exact equilibria. Point convergence, convergence
of empirical play, and convergence in distance to an equilibrium set are
also distinct.

\textbf{Computational representations.} Unless stated otherwise, a complexity
question takes rational data with binary encoding; polynomial time is measured
in total bit length. A fixed-accuracy polynomial algorithm is not necessarily
polynomial in $1/\varepsilon$, and neither is necessarily polynomial in
$\log(1/\varepsilon)$. Succinct games and value, demand, or sampling oracles
must declare their interfaces. Exact real solutions require an explicit
representation; an unspecified list of arbitrary real numbers is not a
finite computational output. A claimed algorithmic barrier is meaningful
only for its specified model and reductions.

\textbf{Dynamic games.} Histories, information, observation, and allowable
randomization are part of the model. A stationary strategy depends only on
the current state; a positional strategy in a deterministic graph game
chooses one action at each controlled vertex. Nash equilibrium at an initial
state and equilibrium after every history are different requirements.
An expectation of a pathwise $\liminf$ payoff, a $\liminf$ of expected averages,
a limiting expected average, and a uniform finite-horizon equilibrium are
different criteria. None is silently substituted for another.

\textbf{Allocation and mechanisms.} A complete allocation assigns every item
exactly once unless otherwise stated. Zero-valued goods, negative values,
capacity constraints, feasibility, lotteries, and copies of goods affect
fairness definitions and are specified locally. Dominant-strategy incentive
compatibility, Bayesian incentive compatibility, universal truthfulness,
and truthfulness in expectation impose different inequalities. Whenever
an objective ratio has zero denominator, its convention is stated or those
instances are excluded.

\textbf{Infinite-dimensional models.} Expectations, integrals, stochastic
equations, and optimization objectives require their local regularity,
integrability, filtrations, and admissible domains. A backward--forward
mean-field system is a boundary-value problem; it is not an initial-value
semiflow without an additional construction. Long-time, large-population,
and vanishing-noise limits require an order of limits and a convergence
topology. If a minimum need not be attained, the objective is its infimum
or supremum and the approximation requirement is stated separately.
% END ECONOMICS STATEMENT
\end{document}
